\documentclass[mathserif, fleqn]{beamer}
% basic text processing/formatting
\usepackage{ucs}                % more unicode
\usepackage[utf8x]{inputenc}
\usepackage[T1]{fontenc}
\usepackage[english]{babel}

% maths stuff
\usepackage{amsmath}
\usepackage{amsfonts}
\usepackage{amssymb}
\usepackage{mathtools}
\usepackage{slashed}            % (Feynman-) slashed symbols
\usepackage{commath}            % for nicer differentials
\usepackage{bm}                 % bold math
\usepackage{dsfont}             % double stroked characters

\usepackage{tikz}               % drawing all manners of things
\usepackage{datetime}           % for custom date and time format
\usepackage{array}              % for matrices in math environment
\usepackage{float}              % place graphics with "H"
\usepackage{braket}             % Dirac notation
\usepackage{placeins}           % FloatBarrier
\usepackage{listings}           % source code listings
\usepackage[xparse, skins]{tcolorbox}   % nice coloured boxes
% for textblock, remove showboxes for release
\usepackage[absolute,overlay,showboxes]{textpos}

%% devel only
% overlay a grid over the slides
\usepackage[texcoord,grid,gridcolor=black!10,subgridcolor=black!5,gridunit=mm]{eso-pic}
\usepackage{lipsum}


%----------------------------------------------------------
% general info
% ISO date format
\newdateformat{isodate}{\THEYEAR-\twodigit{\THEMONTH}-\twodigit{\THEDAY}}

% commands so we can refer to these later
\newcommand{\theauthor}{Jan-Lukas Wynen}
\newcommand{\theinstitute}{Institute for applied \LaTeX}
\newcommand{\thetitle}{Fancy shmancy talk}
\newcommand{\thesubtitle}{Does this subtitle appear anywhere?}
\newcommand{\thedate}{\isodate\today}

\author{\theauthor}
\title{\thetitle}
\date{\thedate}
\institute{\theinstitute}


%----------------------------------------------------------
% setup beamer
\usetheme[titleinfoot, toc]{aiphi}


%----------------------------------------------------------
% tcolorbox
% box for itemizes
\NewTColorBox{itembox}{m}{
  title=#1,
  noparskip,  % not using library parskip
  left=1mm, right=3mm, bottom=3mm, top=3mm,  % less for left to account for itemize
  sharp corners,
  boxrule=0.5mm,
  colframe=aiphiblue,
  colback=white
}

% box with a fancy title
\NewTColorBox{fancybox}{m}{
  enhanced,
  title=#1,
  attach boxed title to top center={
    yshift=-0.25mm-\tcboxedtitleheight/2,  % extra shift is 1/2 boxrule
    yshifttext=1ex-\tcboxedtitleheight/2
  },
  boxed title style={ % draw wedges around the title
    boxrule=0.5mm,
    frame code={ \path[tcb fill frame] ([xshift=-4mm]frame.west)
      -- (frame.north west) -- (frame.north east) -- ([xshift=4mm]frame.east)
      -- (frame.south east) -- (frame.south west) -- cycle; },
    interior code={ \path[tcb fill interior] ([xshift=-2.5mm]interior.west)
      -- (interior.north west) -- (interior.north east)
      -- ([xshift=2.5mm]interior.east) -- (interior.south east) -- (interior.south west)
      -- cycle;}
  },
  noparskip,
  left=3mm, right=3mm, bottom=3mm,
  sharp corners,
  boxrule=0.5mm,
  colframe=aiphiviolet,
  colback=white,
  colbacktitle=aiphiviolet!70!white,
  coltitle=white
}

% box showing a tag besides the title
% arguments: title, tag, tag background, tag foreground
\NewTColorBox{tagbox}{m m O{black!80!white} O{white}}{
  title=#1,
  after title={\hfill\colorbox{#3}{\color{#4}#2}},
  boxsep=0.5mm,
  noparskip,
  left=3mm, right=3mm, bottom=3mm, top=3mm,
  sharp corners,
  boxrule=0.5mm,
  colframe=black!80!white,
  colback=white,
  colbacktitle=white,
  coltitle=black
}

% a color box inside a textblock to allow for arbitrary positioning
\newenvironment{aside}[2]{
  \begin{textblock*}{#1} (#2)
    \begin{tcolorbox}[
      enhanced,
      noparskip,  % not using library parskip
      left=1mm, right=1mm, bottom=1mm, top=1mm,
      sharp corners,
      boxrule=0.01mm,
      colback=aiphiyellow,
      coltext=black!80!white,
      watermark tikz={\draw[line width=2mm] circle (1cm) node{\fontfamily{ptm}\fontseries{b}\fontsize{20mm}{20mm}\selectfont!};},
      ]
    }
    {\end{tcolorbox}\end{textblock*}}

% ----------------------------------------------------------
% TikZ
\usetikzlibrary{tikzmark, calc}


% ----------------------------------------------------------
% code listings
% set typewriter font to courier
\renewcommand{\ttdefault}{pcr}

% Define a new language for C++ to have proper detection of built in datatypes.
% I hope I got all keywords.
\lstdefinelanguage{aiphicpp}{
  morekeywords={class, struct, template, typename, decltype, sizeof, static, if, else, for, while, do, goto, static_cast, dynamic_cast, reinterpret_cast, const_cast, constexpr, const, inline, and, or, xor, bitand, bitor, and_eq, or_eq, xor_eq, return, throw, try, catch, noexcept, namespace, using, typedef, enum, union, beak, continue, switch, case, break, continue, default, delete, new, operator, public, private, protected, explicit, extern, mutable, volatile, override, virtual, final},
  morekeywords=[2]{void, unsigned, int, long, float, double, bool, char, auto},
  morecomment=[l]{//},
  morecomment=[s]{/*}{*/},
  morestring=[b]",
  morestring=[b]'
}

% custom style for code listings
\lstdefinestyle{aiphilst}{
  language=aiphicpp,  % the default language
  basicstyle=\scriptsize\ttfamily\color[RGB]{211,211,211},  % default font size
  breakatwhitespace=false,  % if true, lines are only broken at whitespaces
  breaklines=true,
  showstringspaces=false,  % show marker at spaces in strings
  captionpos=b, % place captions below listin
%
  frame=l,                                  % frame around code, show only to the left
  framerule=1.5pt,                          % frame width
  rulecolor=\color{aiphibordeaux!75!white}, % colour of frame
%
  numbers=left,  % line numbers, one of none, left, right
  numbersep=6pt, % distance of line numbers from code
  numberstyle=\tiny\color{black},  % text style of line numbers
%
  backgroundcolor=\color[HTML]{282b2e},
  commentstyle=\itshape\color[HTML]{8cbbad},
  keywordstyle={[1]\bfseries\color{aiphiblue!75!white}},
  keywordstyle={[2]\color{aiphigreen!75!white}},
  stringstyle=\color{aiphired!70!white},
}
% set this style as default
\lstset{style=aiphilst}


% ----------------------------------------------------------
% handle backup / appendix slides
% place these around backup slides so they are not counted by the progress bar
\newcommand{\backupbegin}{
   \newcounter{framenumberappendix}
   \setcounter{framenumberappendix}{\value{framenumber}}
}
\newcommand{\backupend}{
   \addtocounter{framenumberappendix}{-\value{framenumber}}
   \addtocounter{framenumber}{\value{framenumberappendix}}
}


%----------------------------------------------------------
% abbreviations
\newcommand{\unit}[1]{\,\text{#1}}
\newcommand{\ev}{\,\text{eV}}
\newcommand{\kev}{\,\text{keV}}
\newcommand{\mev}{\,\text{MeV}}
\newcommand{\gev}{\,\text{GeV}}

\renewcommand{\i}{\mathrm{i}}
\renewcommand{\epsilon}{\varepsilon}

\newcommand{\Tr}{\text{Tr}}


% ----------------------------------------------------------
% miscellanious

% remember current item command so we can switch back to it
\let\olditem\item

%%% Local Variables:
%%% mode: latex
%%% TeX-master: "main"
%%% End:


\begin{document}

\begin{frame}
  \titlepage
\end{frame}

\section{A section}
\begin{frame}\frametitle{Overview}
  % different bullets just for this slide
  \setbeamertemplate{itemize item}{}
  \setbeamertemplate{itemize subitem}{\adfrightarrowhead}
  \setbeamertemplate{itemize subsubitem}{--}
  
  \begin{itemize}\setlength{\itemsep}{1.5ex}
  \item \textbf{Motivation}
    \begin{itemize}\setlength{\itemsep}{1ex}
    \item Latex is `fun'
    \item Really nice slides if you put really much effort into it
    \end{itemize}
  \item \textbf{Method}
    \begin{itemize}\setlength{\itemsep}{1ex}
    \item Editor + PDFLatex
    \item Ah, and PDF viewer
      \begin{itemize}
      \item Many choices
      \item Some good some less so
      \item Watch out, sometime coloured areas aren't merged properly!
      \end{itemize}
    \end{itemize}
  \item \textbf{Results}
    \begin{itemize}\setlength{\itemsep}{1ex}
    \item See for yourself
    \item Looks nice, doesn't it?
    \end{itemize}
  \end{itemize}

  \begin{aside}{40mm}{80mm,10mm}
    Note the special itemize markers on this slide
  \end{aside}

  \begin{aside}{30mm}{80mm,65mm}
    And you can place some neat extra notes
  \end{aside}
\end{frame}

\begin{frame}\frametitle{Stuff}
  \begin{itemize}
  \item This is the normal itemize style
    \begin{itemize}
    \item And sub items
      \begin{itemize}
      \item Boring on this level, but don't go here anyway!
      \end{itemize}
    \end{itemize}

  \item But look, nice serifs in equations (optional, maybe you need to distinguish those: $uv$):
    \begin{align*}
      \mathcal{C} = \frac1{Z} \braket{0|\int\,\mathcal{D}\psi\,e^{\i S}|0}
    \end{align*}
  \end{itemize}
\end{frame}

\section{Second}
% we need fragile so the code listing can position itself properly
\begin{frame}[fragile]\frametitle{Code!}
  % language is C++ by default
  \begin{lstlisting}[caption={A test C++ listing}]
// Some nice semi-meta function
template <typename It,
          template <typename...> class F,
          typename... TParam>
struct foreach {
  template <typename... Args>
  static void f(Args && ...args) {
    /* magic! */
    _internal::foreach_impl<typename It::Head,
                            typename It::Tail,
                            F,
                            TParam...>
        ::f(std::forward<Args>(args)...);

    std::cout << "Just testing syntax highlighting\n";
    const char c = 'c';
    std::vector<int> v{3};  // what is the length of this?
  }
};
  \end{lstlisting}
\end{frame}

% fragile again (see above)
\begin{frame}[fragile]\frametitle{\texttt{Python} code}
  % set language of next listing
  \lstset{language=python}
  \begin{lstlisting}[caption={Some Python as well}]
# a hat in a cat
class cat:
  """A cat shaped like an arbitrary hat"""

  def __init__(self, shape):
    self._shape = shape

  def show(self):
    # watch out for auto line breaks:
    print(f"A {self._shape} shaped hat in a now {self._shape} shaped cat")
  \end{lstlisting}
\end{frame}

\begin{frame}\frametitle{Blocks}
  Quite nice but maybe boxes based on \texttt{tcolorbox} are better? See below.
  \begin{block}{A normal block}
    Some normal, boring text, but COLOURS!
  \end{block}
  \begin{exampleblock}{An example}
    This is how you do it.
  \end{exampleblock}
  \begin{alertblock}{Careful now}
    You should really watch out more!
  \end{alertblock}
\end{frame}

\begin{frame}\frametitle{Overlay boxes}
  \begin{textblock*}{60mm} (10mm,20mm)
    \begin{itembox}{first block}
      \begin{itemize}
      \item And-a one
      \item And-a two
      \item And-a
      \end{itemize}
    \end{itembox}
  \end{textblock*}

  \only<2->{
    \begin{textblock*}{60mm} (20mm,28mm)
      \begin{itembox}{second block}
        \begin{itemize}
        \item Foo
        \item Bar
        \item Baz
        \end{itemize}
      \end{itembox}
    \end{textblock*}
  }

  \only<3>{
    \begin{textblock*}{60mm} (30mm,36mm)
      \begin{itembox}{third block}
        \begin{itemize}
        \item En
        \item To
        \item Tre
        \item Fire
        \end{itemize}
      \end{itembox}
    \end{textblock*}
  }
\end{frame}

\begin{frame}\frametitle{Special boxes}
  \begin{fancybox}{Testing...}
    \lipsum[2]
  \end{fancybox}

  \begin{tagbox}{Title}{Tag}%[aiphired]
    Tagged stuff
  \end{tagbox}
\end{frame}

\begin{frame}\frametitle{Arrows}
  % mark start node n1 and target node t1
  Some text with a \tikzmark{n1}note\\[2ex]
  Some more text with a \tikzmark{t1}ref in the middle.
  % draw an arrow from n1 to t1
  \begin{tikzpicture}[remember picture, overlay]
    \draw[->] ([yshift=0.5ex]pic cs:n1) to[bend right] ([yshift=0.5ex]pic cs:t1);
  \end{tikzpicture}

  \vspace{3ex}
  More fancy annotations:
  \begin{align*}
    \ket{\tikzmark{spin1}1, \tikzmark{magn1}m, \tikzmark{mom1}0} = \ldots
  \end{align*}
  % note the manual indent :/
  \hspace{6em}\tikzmark{mom2}Momentum\\
  \hspace{6em}\tikzmark{magn2}Spin projection\\
  \hspace{6em}\tikzmark{spin2}Total spin

  % draw lines and circles
  % all with manual adjustment of positions
  \begin{tikzpicture}[remember picture, overlay]
    \draw let
      \p1 = ([xshift=3pt, yshift=-3pt]pic cs:spin1),
      \p2 = ([xshift=-3pt, yshift=0.7ex]pic cs:spin2)
    in
      (\p1) .. controls (\x1,\y2) .. (\p2)
      node[shape=circle, fill=black, inner sep=1pt] {};

    \draw let
      \p1 = ([xshift=4pt, yshift=-3pt]pic cs:magn1),
      \p2 = ([xshift=-3pt, yshift=0.7ex]pic cs:magn2)
    in
      (\p1) .. controls (\x1,\y2) .. (\p2)
      node[shape=circle, fill=black, inner sep=1pt] {};

    \draw let
      \p1 = ([xshift=3pt, yshift=-3pt]pic cs:mom1),
      \p2 = ([xshift=-3pt, yshift=0.7ex]pic cs:mom2)
    in
      (\p1) .. controls (\x1,\y2) .. (\p2)
      node[shape=circle, fill=black, inner sep=1pt] {};
  \end{tikzpicture}
\end{frame}

\begin{frame}[plain,c]
  % suppress strange item marker in center
  \renewcommand*\item{\olditem[]}
  \begin{center}
    \textcolor{aiphired}{\LARGE Thank you for your attention!}
  \end{center}
\end{frame}

\appendix
\backupbegin

\begin{frame}\frametitle{Appendix 1}
  Some extra stuff\\[2ex]
  Take a look at the progress bar!
\end{frame}

\begin{frame}\frametitle{Appendix 2}
  Even more extra stuff
\end{frame}

\backupend

\end{document}